\documentclass[a4paper,12pt]{article}
\usepackage[utf8]{inputenc}
\usepackage[russian]{babel}
\usepackage{graphicx}
\usepackage{hyperref}
\usepackage{minted}
\usemintedstyle{friendly}
\usepackage{xcolor}
\usepackage{geometry}
\usepackage{amsmath}
\usepackage{enumitem}
\usepackage{float}
\usepackage{indentfirst}
\usepackage{caption}
\usepackage{amsfonts}
\usepackage{tikz}
\usetikzlibrary{shapes.geometric, arrows.meta, positioning}

\documentclass[14pt]{article}
\usepackage[utf8]{inputenc}
\usepackage[russian]{babel}
\usepackage{geometry}
\usepackage{minted}

\geometry{top=2cm, bottom=2cm, left=2.5cm, right=1.5cm}

% Настройки для листингов кода
\setminted{
    fontsize=\small,
    breaklines=true,
    autogobble=true,
    frame=single,
}

\begin{document}

\begin{titlepage}
    \centering
    % Верхний блок
    {\bfseries МИНИСТЕРСТВО НАУКИ И ВЫСШЕГО ОБРАЗОВАНИЯ РФ \\
    ФГБОУ ВО «Белгородский государственный технологический университет им. В.Г. Шухова»\par}
    
    \vfill
    
    % Центральный блок
    {\bfseries\Large Практическое задание №6\par}
    \vspace{0.5cm}
    {\large по дисциплине: Основы программирования и алгоритмизации \par}
    {\large тема: «Линейный поиск» \par}
    
    \vfill
    
    % Блок авторов
    \begin{flushright}
    \begin{tabular}{l}
        \large Выполнил: студент группы ВТ-252 \\
        \large Дронь Дмитрий Алексеевич \\
        \large Проверила: \\
        \large Стеценко Оксана Николаевна
    \end{tabular}
    \end{flushright}
    
    \vfill
    
    % Нижний блок
    {\large Белгород, 2026 г. \par}
\end{titlepage}

\newpage
\begin{center}
\section*{Цель работы}
\end{center}
Цель методических указаний организовать самостоятельную работу студента по выполнению индивидуальных заданий, направленных на закрепление навыков программирования на языке высокого уровня и способов решения математических задач с применением программы для работы с электронными таблицами MS Excel.
\newpage
\begin{center}
\section*{Индивидуальное задание №6}
\end{center}
\textbf{Задача 5.} В Поле размером m * n заполнено прозрачными и непрозрачными
кубиками. Удалить (сделать прозрачными) все непрозрачные кубики, видимые хотя бы с одной из четырёх сторон (видимость анализируется до удаления какого-либо кубика).
\newpage
\section*{Блок-схема}
\begin{center}
\scalebox{0.75}{
\begin{tikzpicture}[
    >=Stealth,
    line/.style={draw,-Stealth,thick},
    block/.style={
        rectangle,
        rounded corners,
        draw,
        align=center,
        text width=3.8cm,
        minimum height=1cm
    },
    blockplain/.style={
        rectangle,
        draw,
        align=center,
        text width=3.8cm,
        minimum height=1cm
    },
    decision/.style={
        diamond,
        draw,
        align=center,
        aspect=2.2,
        text width=3.6cm,
        inner sep=2pt
    },
    io/.style={
        trapezium,
        trapezium left angle=70,
        trapezium right angle=110,
        draw,
        align=center,
        text width=5cm,
        minimum height=1.1cm
    }
]

%------------------ Узлы ------------------

\node[block] (start) at (0,0)
{
input\_matrix(*m, *n, A[MAXN][MAXN])\\
flag = 0\\
tm, tn
};

\node[io] (scanmn) at (0,-3)
{Введите размеры матрицы m n};

\node[decision] (scanfmn) at (0,-6)
{scanf != 2};

\node[blockplain] (flag1) at (-4.85,-8.5)
{flag = 1};

\node[decision] (rangecheck) at (4.5,-8.5)
{tm <= 0 || tn <= 0\\
|| tm > MAXN || tn > MAXN};

\node[blockplain] (flag2) at (-1,-10.5)
{flag = 2};

\node[blockplain] (assignmn) at (10,-10.5)
{
*m = tm;\\
*n = tn;\\
i = 0
};

\node[io] (scanA) at (10,-13)
{Введите поле\\
(\%d строк по \%d чисел)};

\node[decision] (iloop) at (10,-16)
{i < tm \&\& flag == 0};

\node[blockplain] (jinit) at (10,-18.5)
{j = 0};

\node[decision] (jloop) at (10,-21)
{j < tn \&\& flag == 0};

\node[decision] (scanfA) at (10,-24)
{scanf != 1};

\node[blockplain] (jinc) at (10,-26.5)
{j++};

\node[blockplain] (iinc) at (10,-29)
{i++};

\node[block] (ret) at (0,-30)
{возврат flag};

%------------------ Стрелки ------------------

\draw[line] (start) -- (scanmn);
\draw[line] (scanmn) -- (scanfmn);

% scanf != 2 → flag1
\draw[line]
    (scanfmn.west) -- ++(-2.5,0)
    node[midway,above]{да}
    -- (flag1.north);

% scanf OK → rangecheck
\draw[line]
    (scanfmn.east) -- ++(2.15,0)
    node[midway,above]{нет}
    -- (rangecheck.north);

% rangecheck → flag2
\draw[line]
    (rangecheck.west) -- ++(-1.3,0)
    node[midway,above]
    -- (flag2.north);

% rangecheck OK → assignmn
\draw[line]
    (rangecheck.east) -- ++(1.325,0)
    node[midway,above]
    -- (assignmn.north);

\draw[line] (assignmn.south) -- (scanA.north);
\draw[line] (scanA.south) -- (iloop.north);

% i-loop yes → jinit
\draw[line]
    (iloop.south) -- ++(0,-0.3)
    node[midway,right]
    -- (jinit.north);

% jinit → jloop
\draw[line] (jinit.south) -- (jloop.north);

% j-loop yes → scanfA
\draw[line]
    (jloop.south) -- ++(0,-0.3)
    node[midway,right]
    -- (scanfA.north);

% scanfA yes → jinc
\draw[line]
    (scanfA.south) -- ++(0,-0.3)
    node[midway,right]{да}
    -- (jinc.north);

% jinc → jloop (обратная петля)
\draw[line]
    (jinc.east) -- ++(3,0)
    -- ++(0,7.25)
    -- ++(-5.05,0)
    -- (jloop);

% j-loop no → iinc
\draw[line]
    (jloop.west) -- ++(-1,0)
    node[midway,above]
    -- ++(0,-8)
    -- (iinc.west);

% iinc → iloop (обратная петля)
\draw[line]
    (iinc.east) -- ++(4,0)
    -- ++(0,14.8)
    -- ++(-6.05,0)
    -- (iloop);

% i-loop no → ret
\draw[line]
    (iloop.west) -- ++(-3,0)
    node[midway,above]
    -- ++(0,-11.007)
    -- ++(-4,0)
    -- (ret);

% flag1 → ret
\draw[line]
    (flag1.south) -- ++(0,-10)
    -- ++(2,0)
    -- ++(0,-8)
    -- ++(2.85,0)
    -- (ret);

% flag2 → ret
\draw[line]
    (flag2.south) -- ++(0,-8)
    -- ++(-1.85,0)
    -- ++(0,-8)
    -- ++(2.85,0)
    -- (ret);

\end{tikzpicture}

}
\end{center}

\begin{center}
\scalebox{0.55}{
\begin{tikzpicture}[
    >=Stealth,
    line/.style={draw,-Stealth,thick},
    block/.style={
        rectangle,
        rounded corners,
        draw,
        align=center,
        text width=4cm,
        minimum height=1cm
    },
    blockplain/.style={
        rectangle,
        draw,
        align=center,
        text width=3cm,
        minimum height=1cm
    },
    decision/.style={
        diamond,
        draw,
        align=center,
        aspect=2.2,
        text width=2.6cm,
        inner sep=2pt
    }
]

%------------------ Узлы ------------------
\node[block] (input) at (0,0)
{
process\_field( m, n, a[MAXM][MAXN]) \{\\
Result r;\\
r.error\_flag = 0;\\
vis[MAXM][MAXN];\\
i = 0;
};

% === Цикл 1: инициализация vis ===
\node[decision] (icheck1) at (0,-4)
{i < m};

\node[blockplain] (jzero1) at (0,-6.5)
{j = 0};

\node[decision] (jcheck1) at (0,-9)
{j < n};

\node[blockplain] (visset) at (0,-11.5)
{
vis[i][j] = 0;\\
j++;
};

\node[blockplain] (iinc1) at (0,-14.3)
{i++};

% === Между циклами ===
\node[blockplain] (izero2) at (0,-17.5)
{i = 0};

% === Цикл 2: поиск ===
\node[decision] (icheck2) at (0,-20.5)
{i < m};

\node[blockplain] (foundinit) at (0,-23.2)
{
found = 0\\
j = 0
};

\node[decision] (jcheck2) at (0,-26)
{j < n};

\node[decision] (foundcheck) at (0,-28.8)
{found == 0};

\node[decision] (acheck) at (0,-31.8)
{a[i][j] == 1};

\node[blockplain] (visset2) at (0,-34.8)
{
vis[i][j] = 1;\\
found = 1;
};

\node[blockplain] (jinc2) at (0,-37.3)
{j++};

\node[blockplain] (iinc2) at (0,-40)
{i++};

\node[blockplain] (izero3) at (0,-43)
{i = 0};

%------------------ Стрелки ------------------

% input -> icheck1
\draw[line] (input) -- (icheck1);

% icheck1: да -> jzero1
\draw[line]
    (icheck1.south) -- ++(0,-0.2)
    node[pos=0.8,right]
    -- (jzero1.north);

% icheck1: нет -> iinc1 (влево и вниз)
\draw[line]
    (icheck1.west) -- ++(-4,0)
    node[midway,above]
    -- ++(0,-12)
    -- ++(5.85,0)
    -- (izero2.north);

% jzero1 -> jcheck1
\draw[line] (jzero1) -- (jcheck1);

% jcheck1: да -> visset
\draw[line]
    (jcheck1.south) -- ++(0,-0.2)
    node[pos=0.8,right]{да}
    -- (visset.north);
    
% jcheck1: -> iinc1
\draw[line]
    (jcheck1.west) -- ++(-2,0)
    node[pos=0.8,right]
    -- ++(0,-5.3)
    -- (iinc1.west);


% visset -> jcheck1 петля (вправо)
\draw[line]
    (visset.east) -- ++(1.8,0)
    -- ++(0,4)
    -- ++(-3.445,0)
    -- (jcheck1.north);


% iinc1 -> icheck1
\draw[line]
(iinc1.east) -- ++(4,0)
node[pos=0.8,right]
-- ++(0,12)
-- ++(-5.64,0)
-- (icheck1.north);

% izero2 -> icheck2
\draw[line] (izero2) -- (icheck2);

% icheck2: да -> foundinit
\draw[line]
    (icheck2.south) -- ++(0,-0.2)
    node[pos=0.8,right]
    -- (foundinit.north);

% icheck2: нет -> izero3 (вправо и вниз)
\draw[line]
    (icheck2.west) -- ++(-4,0)
    node[midway,above]
    -- ++(0,-21)
    -- ++ (5.85, 0)
    -- (izero3.north);

% foundinit -> jcheck2
\draw[line] (foundinit) -- (jcheck2);

% jcheck2: да -> foundcheck
\draw[line]
    (jcheck2.south) -- ++(0,-0.2)
    node[pos=0.8,right]
    -- (foundcheck.north);

% jcheck2: нет -> iinc2 (влево и вниз)
\draw[line]
    (jcheck2.west) -- ++(-2.5,0)
    node[midway,above]
    -- ++(0,-14)
    -- (iinc2.west);

% foundcheck: да -> acheck
\draw[line]
    (foundcheck.south) -- ++(0,-0.2)
    node[pos=0.8,right]
    -- (acheck.north);


% acheck: да -> visset2
\draw[line]
    (acheck.south) -- ++(0,-0.2)
    node[pos=0.8,right]{да}
    -- (visset2.north);


\draw[line]
    (iinc2.east) -- ++(3.5,0)
    -- ++(0,21)
    -- ++(-5.15,0)
    -- (icheck2.north);

% visset2 -> jinc2
\draw[line] (visset2) -- (jinc2);

% jinc2 -> jcheck2 петля (вправо)
\draw[line]
    (jinc2.east) -- ++(2,0)
    -- ++(0,13)
    -- ++(-3.645,0)
    -- (jcheck2.north);



\end{tikzpicture}
}
\end{center}
\newpage
\section*{Программа}
\begin{minted}{c}
/* Летняя практика. Индивидуальное задание №5 «Линейный поиск»
   Выполнил студент группы ВТ-252 Дронь Д.А.
   Задание 5: В Поле размером m * n заполнено прозрачными и непрозрачными кубиками. Удалить (сделать прозрачными) все непрозрачные кубики, видимые хотя бы с одной из четырёх сторон (видимость анализируется до удаления какого-либо кубика).
*/
#include <stdio.h>

#define MAXM 100
#define MAXN 100

typedef struct {
    int error_flag;
} Result;

// Ввод поля m x n, проверка размеров и значений (0 или 1)
int input_field(int *m, int *n, int a[MAXM][MAXN]) {
    int flag = 0;
    int tm, tn;

    printf("Введите размеры поля m n: ");
    if (scanf("%d %d", &tm, &tn) != 2) {
        flag = 1;
    }

    else if (tm <= 0 || tn <= 0 || tm > MAXM || tn > MAXN) {
            flag = 2;
    }

    else {
        *m = tm;
        *n = tn;
        printf("Введите поле (%d строк по %d чисел: 1 — непрозрачный, 0 — прозрачный):\n", tm, tn);

        int i = 0;
        while (i < tm && flag == 0) {
            int j = 0;
            while (j < tn && flag == 0) {
                if (scanf("%d", &a[i][j]) != 1) {
                    flag = 1;
                } else if (!(a[i][j] == 0 || a[i][j] == 1)) {
                        flag = 3;
                }
                j++;
            }
            i++;
        }
    }

    return flag;
}

/*
Обработка поля:
удалить (сделать прозрачными) все непрозрачные кубики,
видимые хотя бы с одной из 4 сторон
*/
Result process_field(int m, int n, int a[MAXM][MAXN]) {
    Result r;
    r.error_flag = 0;

    int vis[MAXM][MAXN];

    int i = 0;
    while (i < m) {
        int j = 0;
        while (j < n) {
            vis[i][j] = 0;
            j++;
        }
        i++;
    }

    // Видимость слева
    i = 0;
    while (i < m) {
        int found = 0;
        int j = 0;
        while (j < n) {
            if (found == 0) {
                if (a[i][j] == 1) {
                    vis[i][j] = 1;
                    found = 1;
                }
            }
            j++;
        }
        i++;
    }

    // Видимость справа
    i = 0;
    while (i < m) {
        int found = 0;
        int j = n - 1;
        while (j >= 0) {
            if (found == 0) {
                if (a[i][j] == 1) {
                    vis[i][j] = 1;
                    found = 1;
                }
            }
            j--;
        }
        i++;
    }

    // Видимость сверху
    int j = 0;
    while (j < n) {
        int found = 0;
        i = 0;
        while (i < m) {
            if (found == 0) {
                if (a[i][j] == 1) {
                    vis[i][j] = 1;
                    found = 1;
                }
            }
            i++;
        }
        j++;
    }

    // Видимость снизу
    j = 0;
    while (j < n) {
        int found = 0;
        i = m - 1;
        while (i >= 0) {
            if (found == 0) {
                if (a[i][j] == 1) {
                    vis[i][j] = 1;
                    found = 1;
                }
            }
            i--;
        }
        j++;
    }

    // Удаление видимых непрозрачных кубиков
    i = 0;
    while (i < m) {
        j = 0;
        while (j < n) {
            if (a[i][j] == 1 && vis[i][j] == 1) {
                a[i][j] = 0;
            }
            j++;
        }
        i++;
    }

    return r;
}

// Вывод поля
void output_field(int m, int n, int a[MAXM][MAXN]) {
    int i = 0;
    printf("Поле после удаления видимых непрозрачных кубиков:\n");
    while (i < m) {
        int j = 0;
        while (j < n) {
            printf("%d ", a[i][j]);
            j++;
        }
        printf("\n");
        i++;
    }
}

// Основная функция
int main(void) {
    int m, n;
    int a[MAXM][MAXN];

    int err = input_field(&m, &n, a);

    switch (err) {
        case 0:
            break;
        case 1:
            printf("Ошибка: некорректный формат ввода.\n");
            break;
        case 2:
            printf("Ошибка: размеры поля должны быть в диапазоне 1..%d x 1..%d.\n", MAXM, MAXN);
            break;
        case 3:
            printf("Ошибка: элементы поля должны быть только 0 или 1.\n");
            break;
        default:
            printf("Неизвестная ошибка ввода.\n");
            break;
    }

    Result r = process_field(m, n, a);

    switch (r.error_flag) {
        case 0:
            output_field(m, n, a);
            break;
        default:
            printf("Неизвестная ошибка обработки.\n");
            break;
    }

    return 0;
}

\end{minted}
\newpage
\section*{Тесты}
\begin{figure}[H]
    \centering
    \includegraphics[width=0.9\textwidth]{test.png}
    \caption{Результат выполнения программы}
\end{figure}
\newpage
\section*{Вывод}
В ходе выполнения работы мне удалось организовать самостоятельную работу по выполнению индивидуальных заданий, направленных на закрепление навыков программирования на языке высокого уровня и способов решения математических задач с применением программы для работы с электронными таблицами MS Excel.

\end{document}


